\section{Type discipline and foundational scope}
\label{sec:simple-type-theory-03}

Simple types block unrestricted self-application.  If \(x:\alpha\), then
\(xx\) would require simultaneously \(x:\alpha\to\beta\) and \(x:\alpha\),
which the simple grammar cannot satisfy.  This syntactic fact blocks the form
of self-membership used in Russell's paradox.  It does not prove that every
extension of type theory is consistent; consistency remains a metatheorem
about a specified calculus.

Extensionality and the Axiom of Choice at higher types are separate principles.  A system
may take functional extensionality as an axiom,
\[
  \forall f,g{:}(\alpha\to\beta).\,
  \bigl(\forall x{:}\alpha.\,fx=gx\bigr)\to f=g,
\]
while another builds it into equality semantics.  A higher-order Axiom of Choice
operator introduces additional axioms.  Neither principle follows merely from
the type grammar.

\begin{warning}[Proof-assistant metalogics]
A proof assistant based on a dependent type theory, a simple type theory, or a
logical framework does not thereby prove a theorem about an encoded object
logic.  An encoding requires a theorem connecting encoded derivability or
satisfaction to the intended external judgment.
\end{warning}

